% Generated draft: requires mathematical and source review.
\documentclass[11pt]{article}
\usepackage[margin=25mm]{geometry}
\usepackage{fontspec}
\setmainfont{Latin Modern Roman}
\usepackage{amsmath,amssymb,mathtools,unicode-math}
\setmathfont{Latin Modern Math}
\usepackage{enumitem,needspace,xurl,hyperref}
\setlength{\parindent}{0pt}
\setlength{\parskip}{5pt}
\setlength{\emergencystretch}{3em}
\newcommand{\N}{\mathbb{N}}
\newcommand{\Z}{\mathbb{Z}}
\newcommand{\Q}{\mathbb{Q}}
\newcommand{\R}{\mathbb{R}}
\newcommand{\C}{\mathbb{C}}
\newcommand{\F}{\mathbb{F}}
\newcommand{\A}{\mathbb{A}}
\newcommand{\PP}{\mathbb{P}}
\newcommand{\E}{\mathbb{E}}
\newcommand{\eps}{\varepsilon}
\newcommand{\dd}{\,\mathrm{d}}
\begin{document}
\section*{UnsolvedMath ID: 11000014}
\textbf{Source problem code:} catalogue entry
\subsection{Definitions and mathematical statement}
Let $\operatorname{Mod}_g$ be the mapping class group of a closed oriented genus-$g$ surface in the customary convention. A pseudo-Anosov mapping class preserves a pair of transverse measured foliations and scales their measures by reciprocal factors $\lambda,\lambda^{-1}$ with $\lambda>1$. For $\psi\in\operatorname{Mod}_g$, its normal closure is
\[
 \langle\!\langle\psi\rangle\!\rangle_{\operatorname{Mod}_g}
 =\langle h\psi h^{-1}:h\in\operatorname{Mod}_g\rangle.
\]
Question 2.9 asks whether
\[
 \forall\varphi\in\operatorname{Mod}_g\text{ pseudo-Anosov},\quad
 \exists n=n(\varphi)\in\Z_{>0}:\quad
 \langle\!\langle\varphi^n\rangle\!\rangle_{\operatorname{Mod}_g}
 \text{ is a free group}.
\]
No finite rank or uniform exponent is requested. The excerpt does not explicitly fix the genus range or boundary convention.

The following extracted context concerns the earlier Question 2.8, whose statement is absent. It compares Gromov's analogous high-power phenomenon for hyperbolic-type elements of nonelementary word-hyperbolic groups ([Gro], Theorem 5.3.E) with the Margulis normal subgroup theorem for irreducible lattices in real linear semisimple groups without compact factors and of real rank at least two. For $g\geq2$, the source applies this theorem to the symplectic image $\pi(\Gamma)$ of a normal subgroup of $\operatorname{Mod}_g$, and states Proposition 2.10: a normal subgroup containing the Torelli group $\mathcal I_g=\ker\pi$ is commensurable with either $\mathcal I_g$ or $\operatorname{Mod}_g$. Commensurability here means that the intersection has finite index in both groups.

The context further cites Mess's infinitely generated free $\mathcal I_2$, Hain's Malcev Lie algebra $\mathfrak t_g$, and the equivariant Johnson map $\tau:\mathcal I_g\to\bigwedge^3H/H$, where $H=H_1(S_g;\Z)$ and the embedded copy of $H$ is conventionally given by wedging with the symplectic class. It suggests using the invariant Lie-algebra data of $\Gamma\cap\mathcal I_g$ and extending Proposition 2.10 to $K_g=\mathcal I_g(2)$.

Some extracted claims need qualifications: a statement that every finitely generated normal subgroup of $\operatorname{Mod}_2$ has finite-index symplectic image must exclude trivial and finite central cases; and irreducibility arguments about $\bigwedge^3H/H$ require a field, normally $\Q$, and do not imply equality of integral invariant sublattices. These background arguments are not adopted as unqualified theorems.
\subsection{Short English statement}
Does a suitable positive power of every pseudo-Anosov mapping class have free normal closure? The exponent may depend on the mapping class and the free group may have infinite rank.
\subsection{Sources}
\begin{itemize}
\item Local registry: \url{data/unsolvedmath/problems.json}, permanent UnsolvedMath ID $11000014$, exact source code \texttt{catalogue entry}.
\item Farb, \emph{Problems on Mapping Class Groups and Related Topics}, Question 2.9, PDF page 27: \url{https://www.math.uchicago.edu/~farb/papers/mcgbook.pdf}. Supplied access date: 2026-07-29; extraction: pdf-font-structure; page-disambiguated numbering among 329 headings.
\item Supplied reference labels: Gromov [Gro], Theorem 5.3.E; Margulis [Ma]; Mess [Me]; [Fa3]; Hain [Ha3], with a pointer to Section 5.4. Their full bibliography entries and Question 2.8 are not supplied.
\item The registry mentions random-walk/normal-closure progress and analogues involving free groups and their outer automorphisms without precise papers or hypotheses.
\end{itemize}
\end{document}
